% ============================================================
%  Chapitre 3 — Méthodologies : Scrum + CRISP-DM
% ============================================================
\chapter{Méthodologies de Conduite de Projet}
\thispagestyle{fancy}

\section{Introduction}

La réalisation de \farmAI{} v3.0 combine deux méthodologies complémentaires :

\begin{itemize}
  \item \textbf{Scrum} pour la gestion de projet agile, permettant des livraisons
        incrémentales et une adaptation continue aux retours.
  \item \textbf{CRISP-DM} (\textit{Cross-Industry Standard Process for Data Mining})
        pour structurer le cycle de vie des modèles de Data Science intégrés à la
        plateforme.
\end{itemize}

\section{Méthodologie Scrum}

\subsection{Principes Agile}

Le Manifeste Agile \cite{beck2001agile} postule quatre valeurs fondamentales :

\begin{tcolorbox}[infobox, title={Les 4 valeurs Agile appliquées à FARM AI}]
\begin{enumerate}
  \item \textbf{Individus et interactions} $\succ$ processus et outils
        \textit{(équipe réduite, communication directe)}
  \item \textbf{Logiciel fonctionnel} $\succ$ documentation exhaustive
        \textit{(démonstrations sprint par sprint)}
  \item \textbf{Collaboration client} $\succ$ négociation contractuelle
        \textit{(retours continus avec les utilisateurs-cibles agriculteurs)}
  \item \textbf{Réponse au changement} $\succ$ suivi d'un plan
        \textit{(adaptation aux contraintes GPU/cloud découvertes en cours de projet)}
\end{enumerate}
\end{tcolorbox}

\subsection{Équipe Scrum et rôles}

\begin{table}[H]
\centering
\caption{Équipe Scrum et responsabilités}
\label{tab:scrum_team}
\begin{tabular}{llp{7cm}}
\toprule
\textbf{Rôle} & \textbf{Personne} & \textbf{Responsabilités} \\
\midrule
Product Owner  & Mohamed Sayari (+ encadrant) &
  Définition du Product Backlog, priorisation des User Stories,
  validation des incréments fonctionnels \\
Scrum Master   & Mohamed Sayari &
  Organisation des cérémonies, identification et élimination des
  obstacles (blocages cloud, dépendances GPU) \\
Dev Team       & Mohamed Sayari &
  Conception, développement, tests et déploiement de tous les
  composants (backend, frontend, IoT, ML) \\
\bottomrule
\end{tabular}
\end{table}

\subsection{Artefacts Scrum}

\subsubsection{Product Backlog (extrait)}

\begin{longtable}{llp{5.5cm}c}
\caption{Extrait du Product Backlog Smart Farm AI} \label{tab:backlog} \\
\toprule
\textbf{ID} & \textbf{Acteur} & \textbf{User Story} & \textbf{Points} \\
\midrule
\endfirsthead
\toprule
\textbf{ID} & \textbf{Acteur} & \textbf{User Story} & \textbf{Points} \\
\midrule
\endhead
\bottomrule
\endfoot
US01 & Owner  & Créer un compte et me connecter avec JWT + OTP email & 8 \\
US02 & Owner  & Ajouter une ferme avec localisation GPS et superficie & 5 \\
US03 & Owner  & Visualiser les KPIs en temps réel sur le tableau de bord & 13 \\
US04 & Owner  & Recevoir des alertes IoT (SWARM\_RISK, THERMAL\_FAULT) & 8 \\
US05 & Owner  & Scanner une plante malade et obtenir un diagnostic YOLO & 13 \\
US06 & Owner  & Créer un apiaire avec des ruches numérotées (QR codes) & 8 \\
US07 & Owner  & Enregistrer une visite de ruche avec observations & 5 \\
US08 & Owner  & Consulter les graphiques de poids et température des ruches & 8 \\
US09 & Owner  & Dialoguer avec l'assistant en Darija par voix & 13 \\
US10 & Owner  & Envoyer une photo à l'assistant pour analyse multimodale & 8 \\
US11 & Owner  & Voir mes fermes sur une carte et trouver des vétérinaires & 5 \\
US12 & Worker & Scanner une maladie avec la PWA hors réseau & 13 \\
US13 & Worker & Consulter mes tâches assignées et les valider & 5 \\
US14 & System & Détecter les anomalies télémétriques automatiquement & 8 \\
US15 & Admin  & Visualiser les métriques MLflow du dernier entraînement & 5 \\
\bottomrule
\end{longtable}

\subsection{Planification des 6 Sprints}

Le projet a été planifié sur \textbf{14 semaines} réparties en \textbf{6 sprints}
de durée variable selon la complexité des livrables.

\begin{table}[H]
\centering
\caption{Tableau de planification des 6 Sprints Scrum}
\label{tab:sprints}
\small
\begin{tabular}{ccp{2.2cm}p{5.5cm}p{2.5cm}}
\toprule
\textbf{Sprint} & \textbf{Durée} & \textbf{Titre} &
\textbf{Livrables principaux} & \textbf{Jalon} \\
\midrule
S1 & S0--S2\newline(2 sem.) & Infrastructure \& Auth &
  • Preamble FastAPI (routes, middleware)\newline
  • Authentification JWT + OTP email\newline
  • Base de données PostgreSQL (30+ tables)\newline
  • Frontend React/Vite scaffold
  & Auth OK / JWT \\
\midrule
S2 & S2--S4\newline(2 sem.) & IoT \& Télémétrie &
  • Firmware NODE\_A (irrigation ESP32)\newline
  • Firmware NODE\_B (rucher ESP32)\newline
  • Endpoint \api{/api/v1/iot/telemetry}\newline
  • Dashboard télémétrie Recharts\newline
  • Simulation Wokwi validée
  & IoT Live / Dashboard \\
\midrule
S3 & S4--S7\newline(3 sem.) & Module Apiculture &
  • HiveIntel 7 sous-modules\newline
  • CRUD ruchers/ruches/visites\newline
  • Récoltes, stocks, dépenses\newline
  • Planning apicole\newline
  • Intégration télémétrie NODE\_B
  & HiveIntel v1.0 \\
\midrule
S4 & S7--S10\newline(3 sem.) & Vision IA \& Scanner &
  • 8 modèles YOLO v11 intégrés\newline
  • \api{/api/v1/cv/scan} multimodal\newline
  • Scanner Worker PWA (13 modèles)\newline
  • Résultats temps réel + Expert FAB\newline
  • Map Center SIG (Leaflet + OSM)
  & YOLO Prod \\
\midrule
S5 & S10--S12\newline(2 sem.) & Assistant Souverain &
  • Labess-7B via Ollama local\newline
  • RAG ChromaDB (UTAP/AVFA)\newline
  • STT/TTS intégré (WebSpeech API)\newline
  • Vision multimodale LLaVA\newline
  • Lite Mode (Oracle sans GPU)
  & RAG Darija \\
\midrule
S6 & S12--S14\newline(2 sem.) & MLOps \& Finalisation &
  • Pipeline DVC (build\_dataset → train)\newline
  • MLflow tracking (IsolationForest)\newline
  • Gestion ouvriers (CRUD + tâches)\newline
  • Docker Compose + Oracle Cloud\newline
  • Tests finaux + rapport
  & Soutenance v3.0 \\
\bottomrule
\end{tabular}
\end{table}

\subsection{Diagramme Gantt — Planification temporelle}

La figure \ref{fig:gantt_sprints} présente le diagramme Gantt des 6 sprints sur
les 14 semaines du projet.

\begin{figure}[H]
\centering
\begin{tikzpicture}[x=0.9cm, y=1cm]

% Définition des styles locaux
\tikzset{
  sprint/.style={fill=#1, draw=#1!60!black, rounded corners=2pt, opacity=0.85},
  milestone/.style={
    diamond, draw=farmgreen!80!black, fill=farmgreen!30,
    minimum size=0.55cm, inner sep=0pt
  },
  soutenance/.style={
    diamond, draw=red!70!black, fill=red!25,
    minimum size=0.65cm, inner sep=0pt
  }
}

% ── Axe temporel
\draw[->, thick, gray] (-0.3, 0) -- (15.2, 0) node[right, font=\small] {Semaines};

% ── Graduations semaines
\foreach \w in {0,...,14}{
  \draw[gray!40, thin] (\w, 0.05) -- (\w, -0.15);
  \node[below, font=\tiny\color{gray}] at (\w, -0.18) {S\w};
}

% ── Sprints (barres)
\filldraw[sprint=farmblue!50]  (0,  0.15) rectangle (2,  0.75);
\node[font=\scriptsize\bfseries\color{white}, align=center] at (1,   0.45) {S1 -- Infra \& Auth};

\filldraw[sprint=farmblue!65]  (2,  0.15) rectangle (4,  0.75);
\node[font=\scriptsize\bfseries\color{white}, align=center] at (3,   0.45) {S2 -- IoT \& Télémétrie};

\filldraw[sprint=farmblue!75]  (4,  0.15) rectangle (7,  0.75);
\node[font=\scriptsize\bfseries\color{white}, align=center] at (5.5, 0.45) {S3 -- Apiculture};

\filldraw[sprint=farmblue!85]  (7,  0.15) rectangle (10, 0.75);
\node[font=\scriptsize\bfseries\color{white}, align=center] at (8.5, 0.45) {S4 -- Vision IA};

\filldraw[sprint=farmblue!90]  (10, 0.15) rectangle (12, 0.75);
\node[font=\scriptsize\bfseries\color{white}, align=center] at (11,  0.45) {S5 -- Assistant};

\filldraw[sprint=farmblue!95]  (12, 0.15) rectangle (14, 0.75);
\node[font=\scriptsize\bfseries\color{white}, align=center] at (13,  0.45) {S6 -- MLOps};

% ── Jalons (Milestones)
\node[milestone] (m1) at (2, -0.7) {};
\node[below=0.05cm of m1, font=\tiny, align=center] {Auth OK\\JWT};

\node[milestone] (m2) at (4, -0.7) {};
\node[below=0.05cm of m2, font=\tiny, align=center] {IoT Live\\Dashboard};

\node[milestone] (m3) at (7, -0.7) {};
\node[below=0.05cm of m3, font=\tiny, align=center] {HiveIntel\\v1.0};

\node[milestone] (m4) at (10, -0.7) {};
\node[below=0.05cm of m4, font=\tiny, align=center] {YOLO\\Prod};

\node[milestone] (m5) at (12, -0.7) {};
\node[below=0.05cm of m5, font=\tiny, align=center] {RAG\\Darija};

\node[soutenance] (ms) at (14, -0.7) {};
\node[below=0.05cm of ms, font=\tiny\bfseries, align=center, text=red!70!black]
  {Soutenance\\v3.0};

% Lignes de jalons
\foreach \x in {2, 4, 7, 10, 12}
  \draw[dashed, gray!50, thin] (\x, 0) -- (\x, -0.6);
\draw[dashed, red!50, thin] (14, 0) -- (14, -0.6);

% Légende durée totale
\node[right, font=\tiny\color{gray}] at (14.3, 0.45) {14 sem.\\total};

% Légende
\node[milestone, scale=0.8] at (0, -1.55) {};
\node[right, font=\tiny] at (0.3, -1.55) {Jalon sprint};
\node[soutenance, scale=0.8] at (3, -1.55) {};
\node[right, font=\tiny] at (3.3, -1.55) {Soutenance finale};

\end{tikzpicture}
\caption{Diagramme Gantt — Planification Scrum (6 Sprints, 14 semaines)}
\label{fig:gantt_sprints}
\end{figure}

\section{Méthodologie CRISP-DM}

\subsection{Présentation de CRISP-DM}

CRISP-DM (\textit{Cross-Industry Standard Process for Data Mining}) \cite{crisp_dm}
est un processus standard en 6 phases cycliques pour mener des projets de Data Science
de manière structurée et reproductible. Son principal atout est la gestion explicite
des itérations : lorsque l'évaluation révèle des résultats insuffisants, le processus
revient à la phase de modélisation (ou de préparation des données) sans rupture
méthodologique.

\subsection{Application à Smart Farm AI}

La figure \ref{fig:crisp_dm} illustre le cycle CRISP-DM adapté au contexte de FARM AI,
avec les actions concrètes associées à chaque phase.

\begin{figure}[H]
\centering
% ============================================================
%  Diagramme CRISP-DM — Cycle Data Science pour FARM AI
% ============================================================
\begin{tikzpicture}[
  x=1cm, y=1cm,
  phase/.style={
    circle, draw=#1!80!black, fill=#1!25,
    text width=2.1cm, align=center,
    font=\scriptsize\bfseries, inner sep=3pt,
    minimum size=2.3cm
  },
  annotation/.style={
    rectangle, rounded corners=3pt,
    draw=farmgray!60, fill=yellow!15,
    text width=3cm, align=left,
    font=\tiny, inner sep=4pt
  },
  arrow/.style={->, thick, >=Stealth, #1}
]

% ── Nœuds du cycle CRISP-DM (hexagone de rayon 4.2cm) ────────
\def\R{4.2}

% Phase 1 — Compréhension métier (haut)
\node[phase=farmblue] (biz) at (0:\R) {
  Compréhension\\Métier
};

% Phase 2 — Compréhension données (haut-droite)
\node[phase=farmgreen] (data) at (300:\R) {
  Compréhension\\Données
};

% Phase 3 — Préparation données (bas-droite)
\node[phase=farmorange] (prep) at (240:\R) {
  Préparation\\Données
};

% Phase 4 — Modélisation (bas)
\node[phase=farmred] (model) at (180:\R) {
  Modélisation
};

% Phase 5 — Évaluation (bas-gauche)
\node[phase=farmblue!70] (eval) at (120:\R) {
  Évaluation
};

% Phase 6 — Déploiement (haut-gauche)
\node[phase=farmgreen!70] (deploy) at (60:\R) {
  Déploiement
};

% ── Flèches du cycle principal ────────────────────────────────
\draw[arrow=farmblue!60, bend left=20]  (biz)    to (data);
\draw[arrow=farmgreen!60, bend left=20] (data)   to (prep);
\draw[arrow=farmorange!60, bend left=20](prep)   to (model);
\draw[arrow=farmred!60, bend left=20]   (model)  to (eval);
\draw[arrow=farmblue!50, bend left=20]  (eval)   to (deploy);
\draw[arrow=farmgreen!50, bend left=20] (deploy) to (biz);

% ── Flèche d'itération (évaluation → modélisation) ───────────
\draw[arrow=farmred!70, dashed, bend left=40, line width=1.2pt]
  (eval) to node[above, font=\tiny\itshape, pos=0.5]
  {si mAP $<$ 0.70} (model);

% ── Ellipse centrale ──────────────────────────────────────────
\node[
  ellipse, draw=farmgray!40, fill=white,
  text width=2.8cm, align=center,
  font=\small\bfseries\color{farmblue},
  inner sep=6pt
] at (0,0) {
  \farmAI\\[2pt]
  \tiny\color{farmgray}Data Pipeline
};

% ── Annotations FARM AI par phase ─────────────────────────────
% Compréhension Métier → droite
\node[annotation, right=0.6cm of biz] (ann_biz) {
  \textbf{FARM AI :}\\
  • Détection maladies tunisiennes\\
  • Monitoring ruches ESP32\\
  • Anomalies télémétriques\\
  • Assistant Darija
};
\draw[dashed, farmgray!50, thin] (biz.east) -- (ann_biz.west);

% Compréhension Données → bas-droite
\node[annotation, below right=0.3cm and 0.4cm of data] (ann_data) {
  \textbf{Sources :}\\
  • CSV IoT (NODE\_A/B)\\
  • Datasets YOLO Roboflow\\
  • Corpus UTAP/AVFA (RAG)\\
  • Images maladies foliaires
};
\draw[dashed, farmgray!50, thin] (data.south east) -- (ann_data.north west);

% Préparation → bas-gauche
\node[annotation, below left=0.3cm and 0.4cm of prep] (ann_prep) {
  \textbf{DVC Pipeline :}\\
  • build\_dataset (stage 1)\\
  • Normalisation CSV\\
  • Augmentation images\\
  • Split train/val/test
};
\draw[dashed, farmgray!50, thin] (prep.south west) -- (ann_prep.north east);

% Modélisation → gauche
\node[annotation, left=0.6cm of model] (ann_model) {
  \textbf{Modèles :}\\
  • YOLO v11 (8 modèles)\\
  • IsolationForest\\
  • RAG ChromaDB\\
  • Labess-7B LLM
};
\draw[dashed, farmgray!50, thin] (model.west) -- (ann_model.east);

% Évaluation → haut-gauche
\node[annotation, above left=0.3cm and 0.4cm of eval] (ann_eval) {
  \textbf{Métriques :}\\
  • mAP50 / mAP50-95\\
  • Silhouette Score\\
  • Recall / Précision\\
  • BLEU (LLM)
};
\draw[dashed, farmgray!50, thin] (eval.north west) -- (ann_eval.south east);

% Déploiement → haut-droite
\node[annotation, above right=0.3cm and 0.4cm of deploy] (ann_dep) {
  \textbf{Production :}\\
  • MLflow Registry\\
  • Docker + Oracle\\
  • serve\_latest\_model.py\\
  • Alertes temps réel
};
\draw[dashed, farmgray!50, thin] (deploy.north east) -- (ann_dep.south west);

\end{tikzpicture}

\caption{Cycle CRISP-DM appliqué à Smart Farm AI v3.0}
\label{fig:crisp_dm}
\end{figure}

\begin{table}[H]
\centering
\caption{CRISP-DM — Détail des phases pour Smart Farm AI}
\label{tab:crisp_dm_detail}
\begin{tabular}{clp{8.5cm}}
\toprule
\textbf{Phase} & \textbf{Intitulé} & \textbf{Actions dans FARM AI} \\
\midrule
1 & Compréhension métier &
  Identification des pathologies prioritaires (Varroa, mildiou, bruche),
  définition des KPIs (mAP50 $\geq$ 0.85, Silhouette $\geq$ 0.60),
  contraintes opérationnelles (inférence $<$ 200ms) \\
2 & Compréhension données &
  Collecte datasets Roboflow pour YOLO, extraction CSV IoT des nœuds ESP32,
  constitution du corpus RAG (documents UTAP/AVFA en Darija) \\
3 & Préparation données &
  Pipeline DVC \texttt{build\_dataset} : normalisation des CSV, augmentation
  d'images (rotation, bruit), split train/val/test 80/10/10 \\
4 & Modélisation &
  Entraînement 8 modèles YOLO v11, IsolationForest (n=150, contam.=0.05),
  embedding ChromaDB (HNSW cosine), fine-tuning prompts Labess-7B \\
5 & Évaluation &
  mAP50 global 91.4\%, Silhouette Score 0.73 (anomalies),
  \textit{itération si mAP50 $<$ 0.70} (retour phase 4) \\
6 & Déploiement &
  MLflow Registry, \texttt{serve\_latest\_model.py}, Docker Compose,
  Oracle Cloud (LITE\_MODE sans GPU), alertes temps réel \\
\bottomrule
\end{tabular}
\end{table}

\begin{tcolorbox}[technote, title={Critère d'arrêt des itérations CRISP-DM}]
Le projet définit un critère d'arrêt explicite : si le mAP50 moyen des modèles YOLO
est inférieur à \textbf{0.70} après 3 itérations d'entraînement, la phase de
préparation des données est relancée avec une stratégie d'augmentation plus aggressive
(Mosaic×8, MixUp, CutMix) avant de reprendre la modélisation.

\medskip
\textbf{Résultat obtenu :} mAP50 = 0.914 (bee OBB), 0.879 (leaf disease), 0.891
(livestock) — aucune itération supplémentaire requise.
\end{tcolorbox}
